\documentclass[11pt,a4paper]{article}
\usepackage[utf8]{inputenc}
\usepackage[T1]{fontenc}
\usepackage[spanish]{babel}
\usepackage{geometry}
\usepackage{booktabs}
\usepackage{longtable}
\usepackage{array}
\usepackage{hyperref}
\usepackage{float}
\geometry{margin=2.2cm}

\title{Trabajo Practico: Agentes para Catan con Heuristicas y Modelos de Lenguaje}
\author{[Nombre 1] \and [Nombre 2]}
\date{\today}

\begin{document}
\maketitle

\section{Objetivo y alcance}
Este trabajo desarrolla y evalua un agente inteligente para Catan sobre el simulador PyCatan, combinando dos enfoques:
\begin{itemize}
\item un agente heuristico basado en reglas y funciones de evaluacion;
\item una extension con modelos de lenguaje (LLM) para decisiones concretas del juego.
\end{itemize}

El objetivo no es solo maximizar victorias, sino analizar de forma critica la viabilidad real de usar LLM en un entorno multiagente con restricciones de latencia, robustez y coste computacional.

\section{Parte 1: Desarrollo del agente heuristico}
\subsection{Diseno general}
Se implemento el agente \texttt{AlexPelochoJaimeHeuristicAgent} con una politica orientada a:
\begin{itemize}
\item maximizar produccion esperada de recursos (ponderada por probabilidad de dados);
\item mantener diversidad de recursos para evitar cuellos de botella;
\item favorecer expansion y acceso a nodos de alto valor;
\item estabilizar progreso con comercio con banca dirigido por objetivo de construccion.
\end{itemize}

\subsection{Decisiones controladas por el agente}
\begin{enumerate}
\item \textbf{Colocacion inicial} (\texttt{on\_game\_start}): seleccion de nodo y carretera inicial en funcion de puntuacion de asentamiento y potencial de expansion.
\item \textbf{Fase de construccion} (\texttt{on\_build\_phase}): prioridad ciudad $>$ pueblo $>$ carretera $>$ carta de desarrollo.
\item \textbf{Fase de comercio} (\texttt{on\_commerce\_phase}): objetivo dinamico de construccion y seleccion de intercambio con banca para cubrir faltantes.
\item \textbf{Cartas y eventos}: uso de monopolio, year of plenty, road building y movimiento del ladron con reglas de valor esperado.
\end{enumerate}

\subsection{Informacion del estado usada para decidir}
\begin{itemize}
\item recursos disponibles y faltantes por coste de construccion;
\item nodos y carreteras legales;
\item probabilidades de produccion por terrenos adyacentes;
\item puertos y ratios de intercambio;
\item ocupacion propia/rival de nodos para priorizar objetivos.
\end{itemize}

\subsection{Justificacion de las heuristicas}
La combinacion de produccion + diversidad + expansion reduce el riesgo de bloqueos por escasez de recursos y mantiene ritmo de construccion. La prioridad de ciudad y pueblo se alinea con impacto directo en puntos de victoria y economia de largo plazo.

\section{Parte 1.2: Evaluacion mediante benchmarks}
\subsection{Configuracion experimental}
Se han utilizado los benchmarks del proyecto:
\begin{itemize}
\item \texttt{benchmark\_vs\_random.py}
\item \texttt{benchmark\_vs\_agentes\_estandar.py}
\end{itemize}

En esta iteracion se ejecutaron corridas ligeras para validar integracion extremo a extremo con Ollama local y recoger trazas de fallback.

\subsection{Resultados disponibles}
\subsubsection{Benchmark contra random (corrida ligera actual)}
Fuente: \texttt{benchmark\_vs\_random\_resultados.csv}

\begin{table}[H]
\centering
\begin{tabular}{lrrrrrr}
\toprule
Agente & Partidas & Victorias & Win rate & Media puntos & Decisiones LLM & Fallback rate \\
\midrule
Heuristico & 4 & 1 & 0.2500 & 2.00 & 0 & 0.0000 \\
LLM (qwen2.5:0.5b) & 4 & 0 & 0.0000 & 2.00 & 15 & 1.0000 \\
\bottomrule
\end{tabular}
\caption{Resultado actual frente a random.}
\end{table}

\subsubsection{Benchmark contra estandar (corrida ligera actual)}
Fuente: \texttt{benchmark\_vs\_estandar\_resultados.csv}

\begin{table}[H]
\centering
\begin{tabular}{lrrrrrr}
\toprule
Agente & Partidas & Victorias & Win rate & Media puntos & Decisiones LLM & Fallback rate \\
\midrule
Heuristico & 4 & 2 & 0.5000 & 2.25 & 0 & 0.0000 \\
LLM (qwen2.5:0.5b) & 4 & 1 & 0.2500 & 2.25 & 15 & 1.0000 \\
\bottomrule
\end{tabular}
\caption{Resultado actual frente a agentes estandar.}
\end{table}

\subsection{Analisis breve}
\begin{itemize}
\item Ambas corridas se ejecutaron con presupuesto muy reducido (\texttt{max\_rounds=1}), por lo que no deben interpretarse como ranking final de rendimiento.
\item En este escenario rapido, el agente heuristico mantiene mejor rendimiento relativo que la version LLM.
\item El agente LLM registra decisiones, pero en esta configuracion la tasa de fallback es total.
\item Por ello, el rendimiento observado del agente LLM queda fuertemente contaminado por politica heuristica de respaldo.
\end{itemize}

\section{Parte 2: Uso de modelos de lenguaje en el agente}
\subsection{Diseno de la interaccion con el LLM}
La interfaz LLM se diseno para minimizar ambiguedad y tokens:
\begin{itemize}
\item estado de partida serializado en JSON reducido;
\item lista cerrada de acciones legales;
\item salida exigida como JSON con \texttt{action\_index}.
\end{itemize}

Se probaron tres variantes de prompt (\texttt{compact\_json}, \texttt{resource\_focus}, \texttt{safe\_legal}).

\subsection{Integracion en el agente}
El LLM se integra en dos decisiones requeridas por el enunciado:
\begin{enumerate}
\item \texttt{on\_game\_start};
\item \texttt{on\_build\_phase} (decision frecuente).
\end{enumerate}

Si hay error de proveedor, timeout, salida invalida o indice fuera de rango, el sistema aplica fallback a la decision heuristica, garantizando robustez funcional.

\subsection{Modelos y entornos evaluados}
En este momento hay resultados instrumentados con:
\begin{itemize}
\item \textbf{Ollama local}: \texttt{qwen2.5:0.5b}
\end{itemize}

Pendiente para completar al 100\% los requisitos de la practica:
\begin{itemize}
\item al menos un modelo en AWS Bedrock;
\item al menos un modelo via API UPV;
\item comparativa de al menos 3 modelos en total.
\end{itemize}

\subsection{Observaciones tecnicas importantes}
Durante la ejecucion en entorno mixto WSL + PowerShell se observaron dos efectos clave:
\begin{itemize}
\item desde WSL no siempre hay acceso a la API HTTP local de Ollama en Windows;
\item el fallback por CLI puede incrementar latencia de forma significativa y afectar estabilidad de benchmarks largos.
\end{itemize}

Este comportamiento explica parte de las ejecuciones incompletas o muy lentas y el aumento de fallback efectivo.

\section{Comparacion de codigo: Shiyi\_Catan vs PyCatan/Agents}
Se compararon los archivos homonimos en ambos directorios.

\subsection{Resumen}
\begin{itemize}
\item \texttt{AlexPelochoJaimeHeuristicAgent.py}: misma logica funcional.
\item \texttt{AlexPelochoJaimeLLMAgent.py}: misma logica funcional.
\item \texttt{heuristic\_core.py}: misma logica funcional.
\item \texttt{llm\_prompts.py}: mismo contenido funcional.
\item \texttt{llm\_engine.py}: \textbf{si presenta diferencias funcionales relevantes} en \texttt{PyCatan/Agents}.
\end{itemize}

\subsection{Diferencias funcionales relevantes en llm\_engine}
Respecto a \texttt{Shiyi\_Catan/llm\_engine.py}, la version de \texttt{PyCatan/Agents/llm\_engine.py} incluye:
\begin{itemize}
\item fallback por CLI de Ollama cuando falla la API HTTP;
\item resolucion de rutas candidatas para \texttt{ollama.exe} (incluyendo \texttt{CATAN\_OLLAMA\_CLI\_PATH});
\item forzado de respuesta JSON y temperatura 0 en payload de Ollama;
\item parser mas robusto de \texttt{action\_index} (JSON estricto + extraccion relajada por regex).
\end{itemize}

Estas mejoras incrementan tolerancia a fallos de conectividad/formato, pero tambien pueden introducir mayor latencia cuando se activa transporte por CLI.

\section{Uso de herramientas de IA en el desarrollo}
Se utilizaron herramientas de IA generativa (principalmente GitHub Copilot Chat) para:
\begin{itemize}
\item acelerar refactor y estructuracion de agentes;
\item instrumentar benchmark y metricas de fallback;
\item diagnosticar conectividad Ollama y causas de degradacion.
\end{itemize}

Utilidad percibida: alta para iteracion y depuracion.\newline
Limitaciones: necesidad de validacion manual de resultados experimentales y control estricto de reproducibilidad.

\section{Conclusiones}
\begin{enumerate}
\item El agente heuristico presenta un comportamiento consistente y competitivo para el objetivo de la practica.
\item La integracion LLM esta implementada en dos decisiones clave, con mecanismo de seguridad correcto.
\item En las corridas actuales, la tasa de fallback del LLM es elevada, por lo que la evaluacion de calidad intrinseca del modelo todavia es incompleta.
\item Para cierre definitivo se debe completar la matriz de comparacion (prompts y modelos) incluyendo Bedrock y UPV bajo un protocolo experimental uniforme.
\end{enumerate}

\section{Trabajo futuro inmediato}
\begin{itemize}
\item ejecutar benchmarks ligeros comparables por modelo/prompt con control de latencia;
\item anadir tabla de tiempos medios por decision y, cuando sea posible, tokens entrada/salida;
\item consolidar resultados en una figura comparativa final (win rate, fallback rate, latencia);
\item ajustar redaccion final para no superar 10 paginas.
\end{itemize}

\end{document}
